% Source: https://artofproblemsolving.com/community/c3h1643823_number_theory
% Stable identity: aops-topic-1643823-q01
Let $p$ be one of the primes $7$, $13$, $19$. Then $p-1$ is $6$, $12$ or $18$, and each of these divides $36$.

\textbf{Claim: $a^{37}\equiv a\pmod p$.} If $p\mid a$, both sides are $\equiv0$. Otherwise Fermat's little theorem gives $a^{p-1}\equiv1\pmod p$; writing $36=m(p-1)$,
\[
a^{37}=a\cdot\left(a^{p-1}\right)^m\equiv a\cdot1^m=a\pmod p .
\]

\textbf{Combine.} So each of $7$, $13$, $19$ divides $a^{37}-a$. They are distinct primes, so their product $1729$ divides $a^{37}-a$, that is $a^{37}\equiv a\pmod{1729}$. $\blacksquare$

(The same argument shows $a^{k}\equiv a\pmod{1729}$ whenever $36\mid k-1$; this is why $1729$ is a Carmichael number.)
